% =============================================================================
% MPI Parallelisation of Molecular Dynamics Algorithms
% Written Assignment 2 — MPhil Scientific Computing, University of Cambridge
% Candidate Number: XXXXX
% =============================================================================

\documentclass[10pt,twocolumn]{article}

\usepackage[a4paper,
            top=2.0cm,
            bottom=2.3cm,
            left=1.85cm,
            right=1.85cm,
            columnsep=0.62cm]{geometry}
\usepackage{graphicx}
\graphicspath{{../md-mpi/out/plots/}}
\usepackage{amsmath,amssymb}
\usepackage{booktabs}
\usepackage[hidelinks]{hyperref}
\usepackage{siunitx}
\usepackage[numbers,sort&compress]{natbib}
\usepackage{caption}
\usepackage{fancyhdr}
\usepackage{titlesec}
\usepackage{xcolor}
\usepackage{array}
\usepackage{parskip}
\usepackage{dblfloatfix}
\usepackage[section]{placeins}

% ── General spacing / page behaviour ────────────────────────────────────────
\raggedbottom
\setlength{\parindent}{1em}
\setlength{\parskip}{2pt plus 1pt minus 0.5pt}
\setlength{\textfloatsep}{10pt plus 2pt minus 2pt}
\setlength{\floatsep}{8pt plus 2pt minus 2pt}
\setlength{\intextsep}{8pt plus 2pt minus 2pt}
\setlength{\dbltextfloatsep}{10pt plus 2pt minus 2pt}
\setlength{\dblfloatsep}{8pt plus 2pt minus 2pt}
\setlength{\abovecaptionskip}{4pt}
\setlength{\belowcaptionskip}{0pt}

% ── Float placement parameters ─────────────────────────────────────────────
\renewcommand{\topfraction}{0.92}
\renewcommand{\bottomfraction}{0.8}
\renewcommand{\textfraction}{0.08}
\renewcommand{\floatpagefraction}{0.85}
\renewcommand{\dbltopfraction}{0.92}
\renewcommand{\dblfloatpagefraction}{0.85}
\setcounter{topnumber}{2}
\setcounter{bottomnumber}{1}
\setcounter{totalnumber}{3}

% ── Caption style ──────────────────────────────────────────────────────────
\captionsetup{font=small, labelfont=bf, labelsep=period,
              format=plain, justification=justified, singlelinecheck=false}

% ── Section heading styles ─────────────────────────────────────────────────
\titleformat{\section}{\large\bfseries}{\thesection}{0.7em}{}
\titleformat{\subsection}{\normalsize\bfseries}{\thesubsection}{0.6em}{}
\titleformat{\subsubsection}{\normalsize\itshape\bfseries}{\thesubsubsection}{0.5em}{}
\titlespacing*{\section}{0pt}{1.3ex plus 0.3ex minus 0.2ex}{0.5ex plus 0.2ex}
\titlespacing*{\subsection}{0pt}{0.9ex plus 0.2ex minus 0.1ex}{0.35ex plus 0.1ex}
\titlespacing*{\subsubsection}{0pt}{0.7ex plus 0.2ex minus 0.1ex}{0.25ex}

% ── Page style ─────────────────────────────────────────────────────────────
\pagestyle{fancy}
\fancyhf{}
\fancyhead[L]{\small\itshape MPI Parallelisation of Molecular Dynamics}
\fancyhead[R]{\small\itshape Candidate XXXXX}
\fancyfoot[C]{\small\thepage}
\renewcommand{\headrulewidth}{0.4pt}

% ── siunitx config ─────────────────────────────────────────────────────────
\sisetup{separate-uncertainty=true, range-phrase=--}

% =============================================================================
\begin{document}
\hypersetup{pageanchor=false}

% =============================================================================
%  TITLE PAGE  (single column, own page)
% =============================================================================
\onecolumn
\thispagestyle{empty}

\begin{center}
  \vspace*{1.4cm}
  {\LARGE\bfseries MPI Parallelisation of Molecular Dynamics Algorithms\par}
  \vspace{0.7cm}
  \rule{0.62\textwidth}{0.5pt}\\[0.45cm]
  {\large Written Assignment 2\par}
  {\large MPhil in Scientific Computing\par}
  {\large University of Cambridge\par}
  \vspace{0.8cm}
  {\large Candidate Number: \textbf{XXXXX}\par}
  \vspace{0.35cm}
  {\large March 2026\par}
  \vspace{0.55cm}
  {\normalsize Supervisor: Dr Philip Blakely\par}
  {\normalsize Department of Physics\par}
\end{center}

\vfill
\clearpage

% =============================================================================
%  ABSTRACT  (single column, own page)
% =============================================================================
\thispagestyle{empty}

\begin{center}
  \begin{minipage}{0.88\textwidth}
    \setlength{\parindent}{0pt}
    {\bfseries Abstract}\medskip

    A C++17/MPI molecular dynamics solver is implemented and validated
    using three numerical integrators---Forward Euler, fourth-order
    Runge--Kutta (RK4), and Velocity-Verlet---applied to a harmonic oscillator
    verification suite and a Lennard-Jones argon fluid ($N=864$,
    $\rho=0.8\sigma^{-3}$, $T=\SI{94.4}{K}$) following \citealt{Rahman1964}.
    Harmonic oscillator convergence tests confirm global truncation orders of
    $\mathcal{O}(\Delta t)$, $\mathcal{O}(\Delta t^2)$, and
    $\mathcal{O}(\Delta t^4)$ (fitted slopes 1.05, 2.00, 3.94); RK4 error
    saturates at the IEEE~754 double-precision floor (${\approx}3.6\times10^{-15}$).
    For the argon fluid (100-step NVE run), Velocity-Verlet achieves bounded
    energy conservation ($\max|\Delta E/E_0|=0.082\%$) while Forward Euler
    diverges catastrophically ($\max|\Delta E/E_0|=127.6\%$, temperature runaway
    to ${\sim}\SI{400}{K}$), confirming that symplecticity is essential for
    stable MD.
    Radial distribution function peaks at $r/\sigma=1.09$ ($g=2.84$) and
    $r/\sigma=2.07$ ($g=1.25$) compare favourably with Rahman's published
    liquid-structure data.
    Strong scaling at fixed $N=2048$ yields speedup $24.1\times$ at $P=32$
    (efficiency 75\%; Amdahl $f\approx0.010$); size scaling at fixed $P=16$
    gives wall time $\propto N^{1.85}$, consistent with $\mathcal{O}(N^2)$
    direct force evaluation.
  \end{minipage}
\end{center}

\vfill
\clearpage

% =============================================================================
%  TABLE OF CONTENTS  (own page)
% =============================================================================
\thispagestyle{empty}
\setlength{\columnseprule}{0pt}
\tableofcontents
\vfill
\clearpage

% =============================================================================
%  MAIN TEXT  (two column)
% =============================================================================
\twocolumn
\hypersetup{pageanchor=true}
\pagenumbering{arabic}
\setcounter{page}{1}
\thispagestyle{fancy}

% =============================================================================
\section{Introduction}
\label{sec:intro}
% =============================================================================

Molecular dynamics (MD) simulates the time evolution of interacting particle
systems by numerically integrating Newton's equations of motion. Since
\citealt{Rahman1964}'s landmark study of liquid argon, and the earlier
hard-sphere simulations of \citet{AlderWainwright1957}, MD has become central
to computational chemistry, materials science, and biophysics, yielding
thermodynamic properties, structural correlations, and transport coefficients
inaccessible to analytic theory~\citep{FrenkelSmit2002,AllenTildesley2017}.
Applications span protein folding, radiation damage in nuclear materials,
phase-transition kinetics, and the design of pharmaceutical binding sites.

The computational bottleneck is pairwise force evaluation, scaling as
$\mathcal{O}(N^2)$ for direct summation over $N$ particles.
Physically meaningful simulations require $N\sim10^3$--$10^6$ and
$10^6$--$10^9$ timesteps; this is intractable without parallelism.
Distributing the force computation across $P$ cores via the Message Passing
Interface (MPI) reduces the per-rank cost to $\mathcal{O}(N^2/P)$, but
practical efficiency depends on the communication overhead required to
synchronise particle data between ranks.

Modern CPU architectures expose parallelism at multiple levels: vector units
within a core, multi-core sockets, and distributed-memory clusters connected
by a communication fabric. For distributed-memory problems, MPI remains the
standard approach because it gives explicit control over data ownership and
communication, at the cost of making the programmer responsible for where the
data reside and when they are exchanged.

Integrator choice is equally consequential. Long-time stability demands
\emph{symplectic} integrators that preserve the symplectic structure of
Hamiltonian mechanics: non-symplectic schemes exhibit systematic energy drift
regardless of formal truncation order, corrupting thermodynamic averages
over extended simulations. In many-particle systems, exact long-time
trajectory prediction is in any case impossible because of Lyapunov instability;
the practical goal of MD is therefore not to recover one exact future
trajectory, but to sample the correct thermodynamic ensemble and preserve
the correct short-time dynamical behaviour as accurately as possible.

This report presents a C++17/MPI solver implementing Forward Euler, RK4, and
Velocity-Verlet. The integrators are assessed for convergence order,
symplecticity, and energy conservation on the harmonic oscillator and on a
Lennard-Jones (LJ) argon fluid following~\citealt{Rahman1964}. The
particle-decomposition parallelisation is benchmarked for strong scaling
(up to $P=32$) and size scaling (up to $N=2048$), and related to expected
theoretical properties.

% =============================================================================
\section{Methodology}
\label{sec:methodology}
% =============================================================================

\subsection{Numerical Integrators}
\label{sec:integrators}

We consider $m\ddot{\mathbf{x}}_i = \mathbf{F}_i(\{\mathbf{x}\})$,
reformulated as the first-order system
$\mathbf{y}=(\mathbf{x},\mathbf{v})^\top$,
$\dot{\mathbf{y}}=(\mathbf{v},\mathbf{F}/m)^\top$,
with acceleration $\mathbf{a}_n=\mathbf{F}_n/m$.

\subsubsection{Forward Euler}

Plain forward Euler advances position and velocity by simultaneous first-order
Taylor truncations:
\begin{align}
  \mathbf{x}_{n+1} &= \mathbf{x}_n + \mathbf{v}_n\,\Delta t, \label{eq:euler_x}\\
  \mathbf{v}_{n+1} &= \mathbf{v}_n + \mathbf{a}_n\,\Delta t. \label{eq:euler_v}
\end{align}
Both updates are first-order; the global truncation error is
$\mathcal{O}(\Delta t)$. The scheme is not symplectic: it does not preserve
the symplectic two-form of Hamiltonian mechanics, so phase-space volume is
not conserved and total energy drifts monotonically.

\subsubsection{Fourth-Order Runge--Kutta}

RK4 evaluates four stage derivatives of the first-order system:
\begin{align}
  \mathbf{k}_j &= \Delta t\,\mathbf{f}\!\left(\mathbf{y}_n
    + c_j\,\mathbf{k}_{j-1}\right), \quad j=1,\ldots,4, \label{eq:rk4}\\
  \mathbf{y}_{n+1} &= \mathbf{y}_n
    + \tfrac{1}{6}(\mathbf{k}_1+2\mathbf{k}_2+2\mathbf{k}_3+\mathbf{k}_4),
    \nonumber
\end{align}
with coefficients $(c_1,c_2,c_3,c_4)=(0,\tfrac{1}{2},\tfrac{1}{2},1)$.
Global error is $\mathcal{O}(\Delta t^4)$. RK4 is not symplectic.
Crucially, each intermediate stage requires force evaluation at an
intermediate particle configuration; for the LJ system this entails a
separate \texttt{MPI\_Allgatherv} per stage---four per step---making RK4
prohibitively expensive for large-scale interacting systems.

\subsubsection{Velocity-Verlet}

The Velocity-Verlet algorithm~\citep{Verlet1967,Swope1982} implements a
half-kick / drift / force / half-kick sequence:
\begin{align}
  \mathbf{v}_{n+1/2} &= \mathbf{v}_n + \tfrac{\Delta t}{2}\,\mathbf{a}_n,
    \label{eq:vv1}\\
  \mathbf{x}_{n+1} &= \mathbf{x}_n + \Delta t\,\mathbf{v}_{n+1/2},
    \label{eq:vv2}\\
  \mathbf{a}_{n+1} &= \mathbf{F}(\mathbf{x}_{n+1})/m, \label{eq:vv3}\\
  \mathbf{v}_{n+1} &= \mathbf{v}_{n+1/2} + \tfrac{\Delta t}{2}\,\mathbf{a}_{n+1}.
    \label{eq:vv4}
\end{align}
The global truncation error in trajectory coordinates is
$\mathcal{O}(\Delta t^2)$.
The underlying Verlet construction arises from cancellation of odd-order
Taylor terms when combining forward and reverse expansions, giving a local
truncation error $\mathcal{O}(\Delta t^4)$ per step; the
$\mathcal{O}(\Delta t^2)$ global rate reflects accumulation over
$T/\Delta t$ steps and should not be conflated with the local-order result.
The scheme is \emph{symplectic} and time-reversible: it conserves a shadow
Hamiltonian $\tilde{H}=H+\mathcal{O}(\Delta t^2)$, so total energy oscillates
with bounded amplitude indefinitely rather than drifting~\citep{Hairer2006}.
This matters particularly in chaotic many-body systems, where trajectories
diverge exponentially under small perturbations: the correct requirement is
therefore not exact trajectory matching at long times, but accurate ensemble
sampling and faithful preservation of invariants over long runs.
Only one force evaluation per step is required, making Velocity-Verlet the
standard choice for production MD.

\subsection{Lennard-Jones Potential and Force Computation}
\label{sec:lj}

The LJ pair potential and force on particle $i$ are
\begin{align}
  V(r) &= 4\varepsilon\!\left[\left(\frac{\sigma}{r}\right)^{12}
    - \left(\frac{\sigma}{r}\right)^{6}\right], \label{eq:ljpot}\\
  \mathbf{a}_i &= \frac{24\varepsilon}{m}\sum_{j\ne i}
    \frac{\mathbf{x}_i-\mathbf{x}_j}{r_{ij}^2}
    \!\left[2\!\left(\frac{\sigma}{r_{ij}}\right)^{12}
      - \!\left(\frac{\sigma}{r_{ij}}\right)^6\right], \label{eq:ljforce}
\end{align}
with argon parameters $\varepsilon/k_{\rm B}=\SI{120}{K}$,
$\sigma=\SI{3.4}{\text{\AA}}$, $m=66.904\times10^{-27}\,\si{kg}$~\citep{Rahman1964}.
A hard cutoff is applied at $r_c=2.25\sigma$; the geometric constraint
$r_c < L/2\approx5.115\sigma$ (for $N=864$, $L=10.229\sigma$) ensures validity
of the minimum image convention
$\delta x \leftarrow \delta x - L\,\mathrm{round}(\delta x/L)$.
Potential energy is accumulated unconditionally over all $j\ne i$ pairs and
multiplied by 0.5 post-loop (avoiding a branch in the $\mathcal{O}(N^2)$ hot
loop). Newton's third law is not exploited (see \S\ref{sec:parallel}).
The force kernel uses only integer powers: $r^{-2}\to r^{-6}\to r^{-12}$,
with all constants precomputed at compile time.

Initial conditions follow Rahman: an FCC lattice ($N=4k^3$, $k=6$) with
random perturbations $\pm0.01\sigma$, and Maxwell--Boltzmann velocities
via the Box--Muller transform. Temperature is computed from equipartition
using $N_{\rm dof}=3(N-1)$ degrees of freedom after removing centre-of-mass
drift. Box length scales as $L=10.229\sigma\,(N/864)^{1/3}$ to maintain
constant number density.

\subsection{Parallelisation Strategy}
\label{sec:parallel}

\textbf{Particle decomposition.}
Rank $r$ owns a fixed, contiguous block of particles
$[\mathrm{offset}_r,\,\mathrm{offset}_r+N_r^{\rm loc})$, where
$N_r^{\rm loc}=\lfloor N/P\rfloor+(r<N\bmod P\,?\,1:0)$.
Ownership is determined at startup and fixed throughout the simulation.

\textbf{Data layout.}
Particle state is stored in flat interleaved \texttt{std::vector<double>}
arrays of size $3N_r^{\rm loc}$, indexed as \texttt{arr[3*i+d]}.
This layout provides sequential cache access in force loops and maps directly
to \texttt{MPI\_Allgatherv} buffers without packing.
During time stepping, only positions are communicated; velocities are
broadcast once at startup and forces are never communicated.

\textbf{MPI API usage.}
Five collective operations are used:
\begin{enumerate}\itemsep0pt
  \item \texttt{MPI\_Bcast}: rank~0 generates all initial conditions
    (\texttt{std::mt19937\_64}, seed~42) and broadcasts the complete
    $3N$-double position and velocity arrays before the time loop.
  \item \texttt{MPI\_Allgatherv}: the sole collective in the time-stepping
    loop. After updating and wrapping local positions, each rank
    contributes its $3N_r^{\rm loc}$ doubles to reconstruct the global
    position array on every rank, \emph{before} force evaluation.
    Correct sequencing---wrap $\to$ allgather $\to$ force compute---is
    critical: forces require globally consistent positions.
    For the non-interacting harmonic oscillator, this call is bypassed
    entirely.
  \item \texttt{MPI\_Reduce(MPI\_SUM)}: local kinetic and potential
    energies are summed to rank~0 for output.
  \item \texttt{MPI\_Reduce(MPI\_MAX)}: maximum wall time across all
    ranks is reported; separate reduces capture communication statistics.
  \item \texttt{MPI\_Barrier} / \texttt{MPI\_Wtime}: synchronise ranks
    before timing; accumulate \texttt{MPI\_Allgatherv} wall time when
    the timing flag is set.
\end{enumerate}

\textbf{Newton's Third Law trade-off.}
Each rank evaluates forces on its local particles against all $N$ particles,
giving $\mathcal{O}(N^2/P)$ rather than $\mathcal{O}(N^2/2P)$ cost.
This avoids the complex partial-force exchange communication required by
third-law optimisation, keeping a single clean \texttt{MPI\_Allgatherv}
per Euler/Velocity-Verlet step. RK4 requires four gathers per step:
three for intermediate stages and one to refresh forces at the final state.

\textbf{Equilibration and reproducibility.}
A 50-step startup phase applies per-step velocity rescaling
$(\lambda=\sqrt{T_{\rm target}/T_{\rm meas}})$, broadcast via
\texttt{MPI\_Bcast}, followed by a final exact rescale at the
startup-to-production boundary.
The 100-step production run is clean NVE.
All timing runs use identical initial conditions (fixed seed) and are
performed without file output, as required by the assignment brief.

\textbf{Communication-cost perspective.}
The present particle-decomposition algorithm is deliberately simple and
transparent, but it exposes the standard distributed-memory trade-off:
computation is reduced by splitting the force loop across ranks, while
communication incurs latency and bandwidth costs every time global
particle positions are reconstructed. For strong scaling, the amount of
work per rank decreases as $P$ increases, so communication occupies a
larger fraction of the runtime even if its absolute cost remains modest.

\subsection{Validation Methodology}
\label{sec:validation}

Parallel correctness is verified by confirming that energy trajectories
agree between $P=1$ and $P=2$ runs to bitwise precision.
For $P>2$, floating-point non-associativity in \texttt{MPI\_Reduce} tree
reductions introduces differences at ${\sim}10^{-12}$--$10^{-14}$
relative tolerance, consistent with IEEE~754 arithmetic.
This is expected whenever the order of floating-point summation changes
between serial and parallel executions.

% =============================================================================
\section{Results 1: Harmonic Oscillator Verification}
\label{sec:results1}

The harmonic oscillator $\ddot{x}=-\omega^2x$ ($\omega=1$, $x_0=1$,
$\dot{x}_0=0$) has exact solution $x(t)=\cos t$, providing a stringent
benchmark against which convergence and energy behaviour can be assessed
analytically.

\subsubsection*{Trajectories and phase space}

Figure~\ref{fig:ho_traj} shows position, velocity, and phase-space
trajectories at $\Delta t=0.01$ over $T=10$.
All three integrators track the exact solution closely at this resolution;
the right-end segment of the time traces shows that Euler accumulates a
visible phase error near $t=T$.
In phase space (panel~c), Euler traces an outward spiral consistent with
non-symplectic energy growth and violation of Liouville's theorem; both
Velocity-Verlet and RK4 remain on the exact unit circle.

\begin{figure*}[!t]
  \centering
  \includegraphics[width=0.98\textwidth]{results1_figure1abc.png}
  \caption{Harmonic oscillator trajectories at $\Delta t=0.01$, $\omega=1$,
    $T=10$; exact solution in grey.
    \textbf{(a)} Position $x(t)$ and \textbf{(b)} velocity $v(t)$.
    Velocity-Verlet (blue) and RK4 (dark blue dash-dot) overlay the exact
    solution; Euler (orange dashed) shows visible phase deviation near
    $t\approx T$.
    \textbf{(c)} Phase-space portrait.
    Euler spirals outward (non-symplectic); Velocity-Verlet traces a near-exact
    closed orbit (symplectic, Liouville-preserving); RK4 is indistinguishable
    from the exact circle. Central zoom panel in \textbf{(c)} shows the final-sector
    endpoint separation.}
  \label{fig:ho_traj}
\end{figure*}

\subsubsection*{Small vs.\ large timestep}

Figure~\ref{fig:ho_dtcomp} illustrates the effect of timestep size.
At $\Delta t=0.5$, Euler is numerically unstable (amplitude growth to
$|x|\approx8$, spiralling outward in phase space);
Velocity-Verlet remains bounded even at this coarse resolution---a key
property of symplectic integrators, which preserve qualitative dynamics
even when accuracy is poor;
RK4 stays close to the exact trajectory in both time and phase space.

\begin{figure*}[!t]
  \centering
  \includegraphics[width=0.96\textwidth]{results1_figure2.png}
  \caption{Comparison of coarse ($\Delta t=0.5$, dash-dot) and fine ($\Delta t=0.01$,
    solid) timesteps. Left: $x(t)$; right: phase-space portrait.
    Euler (orange, top) is unstable at $\Delta t=0.5$;
    Velocity-Verlet (blue, centre) retains bounded orbits despite coarse resolution;
    RK4 (dark blue, bottom) remains close to exact at both resolutions.}
  \label{fig:ho_dtcomp}
\end{figure*}

\subsubsection*{Convergence order}

Figure~\ref{fig:ho_conv} presents log--log convergence curves for the endpoint
position error $|x_{\rm num}(T{=}10)-x_{\rm exact}(T{=}10)|$ and the RMS
phase-space error, fitted over the asymptotic regime $\Delta t\le0.1$.
Table~\ref{tab:convergence} reports representative errors and fitted slopes.

\begin{table}[!t]
  \centering
  \caption{Global position error $|x_{\rm num}(T{=}10)-x_{\rm exact}|$ and fitted
    convergence slopes.}
  \label{tab:convergence}
  \setlength{\tabcolsep}{4pt}
  \small
  \begin{tabular}{@{}rccc@{}}
    \toprule
    $\Delta t$ & Euler & Vel.-Verlet & RK4 \\
    \midrule
    $10^{0}$ & $7.0\times10^{0}$   & $3.4\times10^{-1}$ & $2.2\times10^{-2}$ \\
    $10^{-1}$& $2.3\times10^{-1}$  & $2.3\times10^{-3}$ & $3.9\times10^{-6}$ \\
    $10^{-2}$& $2.0\times10^{-2}$  & $2.3\times10^{-5}$ & $4.5\times10^{-10}$\\
    $10^{-3}$& $2.0\times10^{-3}$  & $2.3\times10^{-7}$ & $5.3\times10^{-14}$\\
    $5\!\times\!10^{-4}$ & $9.8\times10^{-4}$ & $5.7\times10^{-8}$ &
      $3.6\times10^{-15}$\rlap{$^\dagger$} \\
    \midrule
    Fitted slope & 1.05 & 2.00 & 3.94 \\
    \bottomrule
    \multicolumn{4}{@{}l}{\footnotesize $^\dagger$IEEE~754 double-precision floor.}
  \end{tabular}
\end{table}

Fitted slopes of 1.05, 2.00, and 3.94 confirm the expected first-, second-,
and fourth-order convergence.
At $\Delta t=10^{-2}$, Velocity-Verlet is three orders of magnitude more
accurate than Euler, and RK4 a further five orders more accurate than Verlet.
RK4 saturates at ${\approx}3.6\times10^{-15}$ for $\Delta t=5\times10^{-4}$,
at the IEEE~754 double-precision floor.

The $\mathcal{O}(\Delta t^2)$ global convergence of Velocity-Verlet arises
from accumulation of steps whose local truncation error is formally
$\mathcal{O}(\Delta t^4)$, due to cancellation of odd-order terms in the
symplectic derivation; this local accuracy is not equivalent to fourth-order
global convergence.

\begin{figure*}[!t]
  \centering
  \includegraphics[width=\textwidth]{results1_figure3.png}
  \caption{Convergence to the exact harmonic-oscillator solution ($\omega=1$,
    $x_0=1$, $T=10$).
    \textbf{(a)} Endpoint position error; \textbf{(b)} RMS phase-space error.
    Filled markers: asymptotic fit domain ($\Delta t\le0.1$); open markers: context only.
    Fitted slopes (Table~\ref{tab:convergence}) confirm first-, second-, and fourth-order
    convergence for Euler, Velocity-Verlet, and RK4.
    RK4 saturates at ${\approx}3.6\times10^{-15}$ (IEEE~754 floor).}
  \label{fig:ho_conv}
\end{figure*}

\subsubsection*{Energy conservation}

Figure~\ref{fig:ho_energy} shows the relative energy deviation
$(E-E_0)/|E_0|$ over $T=10$ at $\Delta t=0.01$.
Euler drifts monotonically to $10.5\%$, a direct consequence of its
non-symplectic character.
Velocity-Verlet exhibits bounded oscillation at amplitude ${\sim}2.5\times10^{-5}$
(visible in the inset), consistent with conservation of a shadow Hamiltonian.
RK4's deviation is negligible ($1.4\times10^{-9}\%$) on this short interval,
but its lack of symplecticity means energy would drift on longer timescales.
This directly demonstrates the distinction between short-time local accuracy and
long-time preservation of the invariant structure of the dynamics.

\begin{figure}[!t]
  \centering
  \includegraphics[width=\columnwidth]{results1_figure4.png}
  \caption{Relative energy deviation $(E-E_0)/|E_0|$ for the harmonic oscillator
    ($\Delta t=0.01$, $T=10$).
    Euler (orange) drifts monotonically to $10.5\%$.
    The inset reveals Velocity-Verlet (blue) bounded oscillation at
    ${\sim}2.5\times10^{-5}$ (shadow-Hamiltonian conservation) and negligible
    RK4 deviation ($1.4\times10^{-9}\%$).}
  \label{fig:ho_energy}
\end{figure}

\FloatBarrier

\section{Results 2: Lennard-Jones Argon Fluid}
\label{sec:results2}

The LJ argon system ($N=864$, $\Delta t=\SI{e-14}{s}$) is simulated at
$T=\SI{94.4}{K}$ using the equilibration protocol described in
\S\ref{sec:parallel}, followed by a clean 100-step NVE production run at $P=4$.

\subsubsection*{Energy conservation: mandated 100-step run}

Figure~\ref{fig:lj_energy} shows the signed relative total-energy deviation
$\Delta E/E_0$ for both integrators over the production run.
Velocity-Verlet (panel~a) maintains bounded, oscillatory deviation with
$\max|\Delta E/E_0|=0.082\%$ and mean $0.033\%$, consistent with
shadow-Hamiltonian conservation (cf.~Fig.~\ref{fig:ho_energy}).
Forward Euler (panel~b) diverges catastrophically to $+127.6\%$
within 1~ps, rendering the trajectory completely unphysical.

Three factors contribute to the finite Velocity-Verlet deviations.
The hard cutoff at $r_c$ introduces discontinuities whenever a particle
pair crosses the cutoff radius, causing local energy spikes.
The initial FCC lattice is far from the equilibrium liquid structure,
generating large early-time forces before relaxation.
Discrete-timestep truncation accumulates even in a symplectic scheme.
A shifted-force cutoff---standard in production codes---would eliminate the
force discontinuity at $r_c$.

Exact reproduction of Rahman's energy trajectory is not expected: the
assignment brief explicitly notes sensitivity to initial conditions, and our
FCC starting configuration differs from Rahman's prepared liquid state.
Parallel correctness is confirmed by bitwise agreement of energy trajectories
between $P=1$ and $P=2$ runs.

\begin{figure}[!t]
  \centering
  \includegraphics[width=\columnwidth]{results2_figure6.png}
  \caption{Signed relative energy deviation $\Delta E/E_0$ [\%] during the 100-step
    NVE production run for liquid argon ($N=864$, $\Delta t=\SI{e-14}{s}$,
    equilibrated to $T=\SI{94.4}{K}$).
    \textbf{(a)} Velocity-Verlet: bounded oscillation,
    $\max|\Delta E/E_0|=0.082\%$ (mean $0.033\%$), confirming symplectic
    long-time stability.
    \textbf{(b)} Forward Euler: monotonic drift to $+127.6\%$, demonstrating
    catastrophic non-symplectic energy injection.}
  \label{fig:lj_energy}
\end{figure}

\subsubsection*{Temperature}

Figure~\ref{fig:lj_temp} confirms the temperature evolution.
Velocity-Verlet holds $T=(94.4\pm1.4)\,\si{K}$ throughout, consistent
with NVE energy conservation.
Euler heats from the target $\SI{94.4}{K}$ to ${\sim}\SI{400}{K}$ within 1~ps,
a direct consequence of non-symplectic energy injection; mean temperature
rises to $\SI{185}{K}$ over the 100 steps.
This temperature runaway demonstrates conclusively that Euler is unsuitable
for production MD even from a well-equilibrated initial configuration.

\begin{figure}[!t]
  \centering
  \includegraphics[width=\columnwidth]{results2_figure7.png}
  \caption{Instantaneous temperature during the 100-step argon NVE run
    ($N=864$, $\Delta t=\SI{e-14}{s}$).
    Velocity-Verlet (blue) holds $T\approx\SI{94.4}{K}$ (grey dashed);
    Forward Euler (orange) heats to ${\sim}\SI{400}{K}$, reflecting the energy
    drift of Fig.~\ref{fig:lj_energy}.}
  \label{fig:lj_temp}
\end{figure}

\subsubsection*{Radial distribution function}

To validate the structural properties of the simulated liquid state,
an extended NVE run was performed (20\,000 production steps, Velocity-Verlet,
$P=4$, $N=864$, with a 200-step equilibration discard before $g(r)$
accumulation sampled every 5 steps).
Figure~\ref{fig:lj_rdf} shows $g(r)$ alongside a reference curve derived from
\citealt{Rahman1964} Fig.~2 using WebPlotDigitizer, with paper-anchored points
at identifiable features.

The first peak at $r/\sigma=1.09$ ($g=2.84$), the first minimum at
$r/\sigma=1.55$, the second peak at $r/\sigma=2.07$ ($g=1.25$), and the
long-range tail ($\langle g\rangle=1.002$ for $r>4\sigma$) all compare
favourably with Rahman's published liquid-structure data.
The comparison is qualitative to semi-quantitative: exact agreement is not
expected given different initial configurations and finite run length.
The Rahman reference was constructed from manually anchored digitisation
points and should not be over-interpreted.

\begin{figure*}[!t]
  \centering
  \includegraphics[width=0.86\textwidth]{results2_figure8.png}
  \caption{Radial distribution function $g(r)$ from an extended 20\,000-step
    NVE run (Velocity-Verlet, $N=864$, $P=4$; $g(r)$ sampled every 5 steps
    after a 200-step discard), compared with the reference from
    \citealt{Rahman1964} (orange dashed, digitised).
    First peak: $r/\sigma=1.09$, $g=2.84$; second peak: $r/\sigma=2.07$,
    $g=1.25$; tail $\langle g\rangle=1.002\to1$.
    The comparison is qualitative to semi-quantitative; exact reproduction is
    not expected.}
  \label{fig:lj_rdf}
\end{figure*}

\FloatBarrier

\section{Results 3: Parallel Scaling Analysis}
\label{sec:results3}

Results~3 is a separate timing benchmark, not a repeat of the brief-mandated
100-step argon validation run from Results~2.
All timing runs used Velocity-Verlet, fixed seed 42,
$\Delta t=\SI{e-14}{s}$, and no file output.
Strong scaling used 1000 timed steps at fixed $N=2048$, whereas size scaling
used 1000 timed steps at fixed $P=16$; each configuration was repeated
20 times and the median paired wall/communication sample retained to reduce
measurement noise and fixed-overhead contamination.
All times reported in this section are raw measured seconds from those
benchmark windows, not values normalised to 100 steps.

\subsubsection*{Strong scaling}

Figure~\ref{fig:strong}a--c shows speedup $S(P)=T_1/T_P$, parallel efficiency
$E(P)=S(P)/P$, and the computation/communication breakdown at fixed $N=2048$.
Table~\ref{tab:strong} records the measured values.

\begin{table}[!t]
  \centering
  \caption{Strong scaling results ($N=2048$, Velocity-Verlet, 1000 timed steps;
    medians of 20 no-output repetitions).}
  \label{tab:strong}
  \setlength{\tabcolsep}{3.5pt}
  \small
  \begin{tabular}{@{}rrrrrr@{}}
    \toprule
    $P$ & Wall [s] & Comm [s] & Comm\% & $S(P)$ & $E(P)$ \\
    \midrule
     1 & 45.96 & 0.001 &  0.0 &  1.00 & 1.000 \\
     2 & 23.15 & 0.244 &  1.1 &  1.99 & 0.993 \\
     4 & 12.00 & 0.661 &  5.5 &  3.83 & 0.957 \\
     8 &  6.09 & 0.554 &  9.1 &  7.55 & 0.944 \\
    16 &  3.21 & 0.367 & 11.4 & 14.30 & 0.894 \\
    24 &  2.37 & 0.319 & 13.5 & 19.42 & 0.809 \\
    32 &  1.91 & 0.290 & 15.2 & 24.12 & 0.754 \\
    \bottomrule
  \end{tabular}
\end{table}

Scaling is near-ideal up to $P=8$ (efficiency $>94\%$).
At $P=32$, $S(32)=24.1\times$ with efficiency $75\%$.
Amdahl's law,
\begin{equation}
  S(P)=\frac{1}{f+(1-f)/P}, \label{eq:amdahl}
\end{equation}
fitted via nonlinear least-squares across all $P>1$ values yields
$f\approx0.010$.
This is an \emph{empirical descriptor} of the measured strong-scaling curve
and should not be interpreted as a literal, machine-independent serial
fraction: Amdahl's model assumes constant $f$ and negligible communication,
whereas the \texttt{MPI\_Allgatherv} communication fraction grows monotonically with $P$
(Table~\ref{tab:strong}, rising from $1.1\%$ at $P=2$ to $15.2\%$ at $P=32$),
explaining the visible deviation from the Amdahl fit at high $P$.
This is the expected behaviour of strong scaling: as the work per process
falls, the communication term grows as a fraction of total runtime.
A more realistic extension includes a message-passing overhead,
$T_n = T_s + T_p/N + c\lambda N$, where $\lambda$ represents the cost of
inter-process communication and the final term captures the increasing impact
of communication as process count rises~\citep{Amdahl1967}.

\begin{figure*}[!t]
  \centering
  \includegraphics[width=\textwidth]{results3_figure10abc.png}
  \caption{Strong scaling ($N=2048$ fixed, Velocity-Verlet, 1000 timed steps;
    medians of 20 no-output repetitions).
    \textbf{(a)} Speedup $S(P)$: ideal line (grey dashed), Amdahl fit
    (Eq.~\ref{eq:amdahl}, $f=0.010$, orange).
    Near-ideal to $P=8$; $S(32)=24.1\times$.
    \textbf{(b)} Parallel efficiency $E(P)$: above $94\%$ for $P\le8$,
    $75\%$ at $P=32$.
    \textbf{(c)} Critical-path communication (orange) vs.\ remaining runtime
    (blue); communication fraction grows from $0\%$ to $15\%$,
    explaining the Amdahl fit deviation at high $P$.}
  \label{fig:strong}
\end{figure*}

\subsubsection*{Size scaling}

Figure~\ref{fig:size} shows wall time and its components as a function of $N$
at fixed $P=16$, ranging from $N=108$ to $N=2048$.
Empirical log--log fits give wall time $\propto N^{1.85}$ and compute-only
$\propto N^{1.93}$, close to the expected $\mathcal{O}(N^2)$ for direct
all-pairs force evaluation.
These are fits over a limited range of $N$ and should be treated as such.
The sub-quadratic wall-time exponent reflects the $\mathcal{O}(N)$
\texttt{MPI\_Allgatherv} communication, which contributes a sub-quadratic
component; the compute-only exponent converges toward~2 as $N$ grows,
confirming the $\mathcal{O}(N^2)$ force loop asymptotically dominates.
Communication fraction falls monotonically from $51\%$ at $N=108$ to $13\%$
at $N=2048$ (Fig.~\ref{fig:size}b), confirming that $\mathcal{O}(N^2)$ force
cost grows faster than $\mathcal{O}(N)$ communication volume.

Although only strong and size scaling are required here, it is helpful to note
the complementary weak-scaling viewpoint: if work per process is held roughly
constant while both problem size and process count increase together, one
expects communication effects to be isolated more clearly, rather than being
amplified by shrinking local workload.

\begin{figure*}[!t]
  \centering
  \includegraphics[width=\textwidth]{results3_figure9ab.png}
  \caption{Size scaling ($P=16$ fixed, Velocity-Verlet, 1000 timed steps;
    medians of 20 no-output repetitions).
    \textbf{(a)} Log-log wall time (blue circles), compute-only (blue triangles),
    and communication time (orange squares) vs.\ $N$, with $N^2$ reference.
    Empirical fits: $T_{\rm wall}\propto N^{1.85}$, $T_{\rm compute}\propto N^{1.93}$.
    \textbf{(b)} Communication fraction vs.\ $N$: monotonic decrease from $51\%$
    to $13\%$ confirms $\mathcal{O}(N^2)$ force cost grows faster than
    $\mathcal{O}(N)$ \texttt{MPI\_Allgatherv} volume.}
  \label{fig:size}
\end{figure*}

\FloatBarrier

% =============================================================================
\section{Conclusions}
\label{sec:conclusions}
% =============================================================================

The harmonic oscillator tests confirm the expected convergence orders (1.05,
2.00, 3.94 for Euler, Velocity-Verlet, and RK4) and the decisive role of
symplecticity in long-time energy conservation.
Velocity-Verlet's bounded energy oscillation ($0.002\%$) versus Euler's
monotonic $10.5\%$ drift shows that preserving the underlying Hamiltonian
structure matters more for long-time MD stability than formal local order alone.

For the LJ argon fluid, Velocity-Verlet conserves energy to $0.082\%$ over
the mandated 100-step NVE run, while Euler diverges to $128\%$ with
temperature runaway to $\sim\SI{400}{K}$.
This demonstrates conclusively that Euler is unsuitable for production MD.
RK4 achieves the highest formal accuracy (down to the IEEE~754 floor at
$3.6\times10^{-15}$) but requires four force evaluations and four
\texttt{MPI\_Allgatherv} calls per step; it is also non-symplectic and would
drift on long timescales. Velocity-Verlet, requiring one force evaluation
and one collective per step, is the correct choice.

Strong scaling achieves $24.1\times$ speedup at $P=32$ with $75\%$
efficiency (empirical Amdahl $f\approx0.010$), limited at high $P$ by
growing \texttt{MPI\_Allgatherv} overhead not captured by the simple
Amdahl model.
The fundamental bottleneck is $\mathcal{O}(N)$ communication volume per step:
the full position array must be gathered every timestep. In distributed-memory
computing, this cost is ultimately controlled by latency and bandwidth as well
as algorithmic structure, so parallel performance is never determined by
flop-count alone.

Production codes such as LAMMPS~\citep{Thompson2022} and GROMACS use
\emph{spatial domain decomposition} with \emph{neighbour lists} (Verlet or
cell lists) to reduce force evaluation to $\mathcal{O}(N)$ and limit
communication to $\mathcal{O}(N^{2/3}/P)$ halo exchange---a fundamental
scaling advantage.
Additional improvements include shifted-force cutoffs to eliminate energy
discontinuities at $r_c$, and Ewald summation for long-range electrostatics.
The complementary perspective of Gustafson--Barsis's law is relevant here:
the size-scaling results show the code can tackle larger problems on more
cores, even if the strong-scaling efficiency is not perfect~\citep{Gustafson1988}.

\section*{Acknowledgements}
The author thanks Dr Philip Blakely for supervision and the Cambridge MPhil
Scientific Computing programme for compute access.

% =============================================================================
% REFERENCES
% =============================================================================
\bibliographystyle{plainnat}

\begin{thebibliography}{99}

\bibitem[Alder \& Wainwright(1957)]{AlderWainwright1957}
Alder B.~J., Wainwright T.~E., 1957, J.\ Chem.\ Phys., 27, 1208

\bibitem[Allen \& Tildesley(2017)]{AllenTildesley2017}
Allen M.~P., Tildesley D.~J., 2017, Computer Simulation of Liquids, 2nd edn.
Oxford University Press

\bibitem[Amdahl(1967)]{Amdahl1967}
Amdahl G.~M., 1967, in \emph{Proceedings of the AFIPS Spring Joint Computer Conference}, pp. 483--485

\bibitem[Frenkel \& Smit(2002)]{FrenkelSmit2002}
Frenkel D., Smit B., 2002, Understanding Molecular Simulation, 2nd edn.
Academic Press

\bibitem[Gustafson(1988)]{Gustafson1988}
Gustafson J.~L., 1988, Commun.\ ACM, 31, 532

\bibitem[Hairer et~al.(2006)]{Hairer2006}
Hairer E., Lubich C., Wanner G., 2006, Geometric Numerical Integration:
Structure-Preserving Algorithms for Ordinary Differential Equations, 2nd edn.
Springer

\bibitem[Rahman(1964)]{Rahman1964}
Rahman A., 1964, Phys.\ Rev., 136, A405

\bibitem[Swope et~al.(1982)]{Swope1982}
Swope W.~C., Andersen H.~C., Berens P.~H., Wilson K.~R., 1982,
J.\ Chem.\ Phys., 76, 637

\bibitem[Thompson et~al.(2022)]{Thompson2022}
Thompson A.~P. et al., 2022, Comput.\ Phys.\ Commun., 271, 108171

\bibitem[Verlet(1967)]{Verlet1967}
Verlet L., 1967, Phys.\ Rev., 159, 98

\end{thebibliography}

\end{document}
